\section{The translated origin and integer lattice points}
\label{sec:area_law_dyadic_origin}

The geometric construction in \cite[Section~11]{OpenAI2026AreaLaw}
uses a translated origin whose coordinates, sum and difference are
nonintegral.

% Source: 10-geometry.tex, prop:two-families, lines 150--154 and 545--559,
% revision adc7f1241b42e322a6451854ab7e4b4c146bf78a.
% Actual edge containment in these four line families is a separate result.

\begin{definition}[The prescribed origin]
    \label{def:area_law_dyadic_origin}
    \lean{TNLean.PEPS.AreaLaw.Geometry.dyadicOrigin}
    \leanok
    Set $o=(\sqrt 2,\sqrt 3)\in\mathbb R^2$.
\end{definition}

\begin{theorem}[Four nonintegrality assertions]
    \label{thm:area_law_dyadic_origin_nonintegral}
    \lean{TNLean.PEPS.AreaLaw.Geometry.dyadicOrigin_nonintegral}
    \leanok
    \uses{def:area_law_dyadic_origin}
    None of $o_1$, $o_2$, $o_1+o_2$ and $o_1-o_2$ is an integer.
\end{theorem}
\begin{proof}\leanok
    \uses{def:area_law_dyadic_origin}
    Squaring positive rational bounds gives
    $4/3<\sqrt 2<2$ and $5/3<\sqrt 3<2$; also
    $\sqrt 2<\sqrt 3$. Thus $o_1,o_2\in(1,2)$,
    $o_1+o_2\in(3,4)$ and $o_1-o_2\in(-1,0)$.
    No integer lies strictly between consecutive integers.
\end{proof}

\begin{theorem}[Integer-translated supporting lines avoid the lattice]
    \label{thm:area_law_dyadic_origin_lines}
    \lean{TNLean.PEPS.AreaLaw.Geometry.dyadicOrigin_supporting_lines_avoid_lattice}
    \leanok
    \uses{def:area_law_dyadic_origin}
    For every $m\in\mathbb Z$, none of the four lines
    \begin{align}
        x_1     & =o_1+m,     \\
        x_2     & =o_2+m,     \\
        x_1+x_2 & =o_1+o_2+m, \\
        x_1-x_2 & =o_1-o_2+m
    \end{align}
    contains a point of $\mathbb Z^2$.
\end{theorem}
\begin{proof}\leanok
    \uses{thm:area_law_dyadic_origin_nonintegral}
    If $x\in\mathbb Z^2$ satisfies one of these equations, subtracting
    $m$ expresses the corresponding coordinate, sum or difference of
    $o$ as an integer. Each possibility contradicts nonintegrality.
\end{proof}
